% Standalone section fragment; needs amsmath, amssymb, amsthm.
\section{A hybrid controller with a new proof interface}
\label{sec:hybrid-proof}

The new machines combine two previously supplied operations whose
combination is not one of the supplied command interfaces.  Modus ponens
returns its conclusion in the visible top, preserves the offset containing
that conclusion, and pushes no dummy cell.  Shift tests only the cached
head.  PA194H realizes these commands directly; PA194HS uses the supplied
symmetric comparison optimization to implement exactly the same commands
with three fewer primitive sites.  The new contribution is the combination,
its exact proof compiler and its literal 194-state implementations.  Neither
the no-push idea nor the symmetric comparison is claimed as new in isolation.

\subsection{Numerical commands}

Write $p(a,b)=(a+b)(a+b+1)/2+b$, $[a;w]=p(a,w)$,
$I(a,b)=4p(a,b)$ and $U(a,B)=4p(B,a)+2$.  The total projections satisfy
$\operatorname{hd}[a;w]=a$, $\operatorname{tl}[a;w]=w$, and
$\operatorname{hd}(0)=\operatorname{tl}(0)=0$.  In particular $[0;0]=0$.
No stack-length argument below assumes otherwise.

A command store is $(t,w,d,h)$.  Its two changed commands are
\[
 S(t,w,d,h)=
 \begin{cases}
 \mathsf{halt},&h=1,\\
 (h,\operatorname{tl}w,d,\operatorname{hd}w),&h\ne1,
 \end{cases}
\]
and, writing $b=d+x$,
\[
 \mathsf{MP}_x(t,w,d,h)=
 \begin{cases}
 (b,w,b,0),&I(t,b)=h,\\
 (0,w,b,0),&I(t,b)\ne h.
 \end{cases}
\]
AX is unchanged from the separate-diagonal-head calculus.  It pushes ten
entries, including the saved old top.  Put $D=d+x$, $H=h+y$, and $A=c$.
In stack-top order its nine generated entries are
\[
 (q_1,\ldots,q_9)=
 (\Gamma(A),L_6,L_5,L_4,L_3,L_2,L_1,\rho_8(D,A,H,d),U(d,h)).
\]
Here $\Gamma(B)=I(B,U(B,B))$, the six $L_i$ are the unchanged universal
heads, $\rho_8$ is the unchanged dominated implicational head, and the
visible output is the unchanged guarded separation gate $K(H,D,A)$.
Thus
\[
 \mathsf{AX}_{c,x,y}(t,w,d,h)=
 (K(H,D,A),[q_1,\ldots,q_9,t;w],0,0).
\]
All $q_1,\ldots,q_8$ and the gate have residue zero modulo four;
$q_9$ has residue two.  Hence none equals one, including the failed-guard
case where $K=0$.  The executable free-term check verifies the entire
emitter stack and top, not merely this residue observation.

\subsection{Uniform word identities}

Words execute from left to right.  For every incoming store,
\begin{align*}
 \mathsf{AX}_{c,x,y};S^9(t,w,d,h)
   &=(\rho_8(D,c,H,d),[t;w],0,U(d,h)),\\
 \mathsf{AX}_{c,x,y};S^{10}(t,w,d,h)
   &=(U(d,h),w,0,t).
\end{align*}
These executions never halt: they test only the cleared initial cache and
the nine generated heads, not the saved arbitrary $t$.

Define
\[
 \mathsf{Push}=\mathsf{AX}_{0,0,0};S^9;\mathsf{MP}_0,
 \quad \mathsf{Pop}=\mathsf{AX}_{0,0,0};S^{11},
 \quad \mathsf{Discard}=\mathsf{MP}_0.
\]
The exact contracts are
\begin{align*}
 \mathsf{Push}(t,w,d,h)&=(0,[t;w],0,0),\\
 \mathsf{Pop}(t,[a;w],d,h)&=
   \begin{cases}\mathsf{halt},&t=1,\\(t,w,0,a),&t\ne1,
   \end{cases}\\
 \mathsf{Discard}(t,w,0,h)&=(0,w,0,0).
\end{align*}
For Push, the final MP compares a residue-zero implication with the
residue-two universal head and therefore fails.  Pop's eleventh Shift
tests precisely the saved top, and is its only possible halt.  For
Discard, the candidate conclusion is zero; both the success and failure
branches return the same store.  Consequently Discard clears an arbitrary
old cache without any equality side condition.  This is shorter than the
multi-command discard required by the supplied cached-result controller.

The empty-stack alias is harmless: every displayed equality is an equality
of the total numerical projections.  In particular Push sends an all-zero
store to an all-zero store; it does not create a distinguishable zero cell.

Let $\mathsf{Offset}_i=\mathsf{MP}_i$.  From $(0,w,0,0)$ it returns
$(0,w,i,0)$ for every $i$, including zero.  Generalization is compiled as
\[
 \mathsf{Gen}_i=
 \mathsf{Push};\mathsf{Offset}_i;S;\mathsf{AX}_{i,0,0};S^{10}.
\]
It sends $(a,w,0,0)$ to $(U(i,a),w,0,0)$, even for $a=1$.
The sole Shift before AX tests zero, then places $a$ in the cache; AX
consumes that cache before any further Shift.

For $1\le j\le7$, the ordinary head selector
\[
 \mathsf{Discard};\mathsf{AX}_{A,D,H};S^{j+1};\mathsf{Pop}^{9-j}
\]
sends arbitrary $(t,w,0,0)$ to $(q_j,w,0,0)$.  The selected head has
residue zero and so every Pop is safe.  The last consumed stack entry is
the saved zero.  The diagonal head is $j=1$, giving explicitly
$\mathsf{Discard};\mathsf{AX}_{B,0,0};S^2;\mathsf{Pop}^{8}$ for
$\Gamma(B)$.  The gate selector is
\[
 \mathsf{Discard};\mathsf{AX}_{c,b,a};\mathsf{Pop}^{10}.
\]
For $u\le a$, the dominated implicational selector is
\[
 \mathsf{Discard};\mathsf{Offset}_u;
 \mathsf{AX}_{b,a-u,c};S^9;\mathsf{Pop}.
\]
It returns $\rho_8(a,b,c,u)$ with the original stack and cleared caches.

\subsection{The new MP suffix and proof compiler}

Define
\[
 \mathsf{MPSuffix}_b=S;S;\mathsf{MP}_b;\mathsf{Pop}.
\]
For arbitrary $a,b,w$ it obeys the exact alternative
\[
 \mathsf{MPSuffix}_b(0,[a,I(a,b),0;w],0,0)=
 \begin{cases}
 \mathsf{halt},&a=1\text{ or }b=1,\\
 (b,w,0,0),&a\ne1\text{ and }b\ne1.
 \end{cases}
\]
The first Shift loads $a$ into the cache.  The second either halts or
establishes $(a,[0;w],0,I(a,b))$.  MP succeeds, leaving
$(b,[0;w],b,0)$.  The final Pop either detects $b=1$, or consumes the
reserved zero and clears the offset without changing $b$ or the old
stack.  Reserving this zero is necessary: without it, Pop would consume
the first entry of the arbitrary underlying stack.  At $a=b=w=0$ the
same calculation applies, because $I(0,0)=0$ and $[0;0]=0$.

Compile each finite numerical proof tree $p$ into a word $W_p$.  Use the
selectors above for axioms and $W_{p_0};\mathsf{Gen}_i$ for generalization.
At an MP node with minor premise $a$, major premise $I(a,b)$, and
subproofs $p_0,p_1$ respectively, set
\[
 W_p=\mathsf{Discard};\mathsf{Push};
       W_{p_1};\mathsf{Push};W_{p_0};\mathsf{Push};\mathsf{MPSuffix}_b.
\]
Induction on the proof tree proves that for every $t,w$, execution from
$(t,w,0,0)$ either halts early or returns
$(\operatorname{concl}(p),w,0,0)$.  The first two commands reserve the
dummy zero.  The next two compiled proofs and Pushes establish the exact
suffix precondition.  Generalization tolerates a premise equal to one,
and every earlier halt remains a halt under word extension.

If $p$ proves one, then
\[
 \mathsf{AX}_{0,0,0};W_p;\mathsf{Pop}
\]
halts from \emph{every} incoming store.  The preliminary AX clears $d,h$;
the compilation invariant tolerates its arbitrary resulting top and stack.
If no earlier halt occurs, the final Pop tests the proven top one before
it depends on any value of the underlying stack.  This dirty-store
property is the precise interface needed by the unchanged exhaustive
command-word scheduler, whose state is carried between trials.

For completeness, the scheduling argument is finite and cofinal, rather
than an assertion that every trial starts with an empty store.  The
inherited command encoding has
\[
 B_S(z)=p(0,z),\qquad B_{\mathsf{MP}_x}(z)=p(1,p(x,z)),\qquad
 B_{\mathsf{AX}_{c,x,y}}(z)=p(c+2,p(y,p(x,z))).
\]
For continuation words define $E(\epsilon)=0$ and
$E(c::w)=1+B_c(E(w))$.  Put
$z=E(\mathsf{AX}_{0,0,0};W_p;\mathsf{Pop})$.
Prefixing that universal halting word by $\mathsf{AX}_{k,0,0}$ gives
the fresh-candidate codes
\[
                 n_k=p(k+2,p(0,p(0,z))).
\]
They are strictly increasing and unbounded in $k$.  Their prefix command
can change the store arbitrarily, because the following word halts from
every store.  The fresh candidate counter starts at two, so omitting zero
and one loses none of this cofinal family.  For any positive continuation
$l$, the scheduler decodes $l-1$ and its next residual continuation is at
most $l-1$.  Every finite trial therefore terminates or halts; the
unchanged arithmetic source loops also have primitive-recursive terminating
bounds.  Bounded induction through the earlier candidate numbers reaches
one chosen $n_k$, unless a halt occurs first.  The continuation distinction
is essential: $E(w)$ is a positive continuation code, while the first
candidate is its outer command code without the added one.  PA proves
the coding identities and finite execution induction directly.

\subsection{Soundness and PA-internal calibration}

Maintain a numerical-calculus derivation for the visible top, cached head
and every live entry in the canonical finite Cantor stack.  Code zero has
the fixed derivation $\rho_8(0,0,0,0)=0$; it supplies the virtual zero tail.
No derivability invariant is imposed on the offset $d$.

AX preserves this annotation by the unchanged axiom instances and theorem
generalization.  MP success obtains $b$ by modus ponens from the annotated
minor $t$ and major $I(t,b)=h$; MP failure returns the fixed truth code.
Shift merely moves annotated entries and halts only on an annotated proof
of one in the cache.  The annotation updates and the compiler above are
primitive recursive operations on finite objects.  Ordinary induction in
PA verifies their correctness.  No reflection or soundness axiom for PA
is used.

The unchanged separate-head logical calibration supplies PA-verifiable
proof translations between refutations of PA and proofs of numerical
code one.  Its original mathematical texts are retained in
\texttt{references/logical/}; the supplied project's missing end-to-end
Lean/iPRA reductions are not treated as completed.  Combining those
inherited translations with the new command proof and unchanged exhaustive
scheduler proves, in PA,
\[
 \operatorname{Halt}(R_H)\ \longleftrightarrow\
 \exists p\;\operatorname{Prf}_{\mathrm{PA}}(p,\ulcorner0=1\urcorner).
\]
This is an ordinary mathematical proof of a PA-provability statement,
not a newly compiled proof-assistant theorem.

\subsection{Exact source binding and the symmetric implementation}

The source has nine logical registers $(N,T,W,A,D,H,X,Y,L)$.
The new MP subtree replaces \texttt{node163=(node6,node55)} with
\texttt{(node6,node54)}.  As in PA194N, it removes the dummy-zero push;
on success the retained seven-site copy-and-restore tail sends
$(T,D,H,X)=(0,b,0,0)$ to $(b,b,0,0)$ in $7b+2$ primitive steps.

The Shift source in PA194H first clears the old top, copies $H$ into $T$
and $X$, restores $X$ to $H$, increments $X$ and jumps to the old root
comparison.  From arbitrary $t,h$ with $X=0$, it reaches
$(T,H,X)=(h,h,1)$.  The unchanged destructive equality guard halts just
when $h=1$; otherwise it clears both compared registers and continues to
the unchanged Cantor decoder.  This is exactly the displayed cache-only
Shift command.

PA194HS instead omits the restore loop and increments $H$, arriving at
$(T,H,X)=(h,1,h)$.  Equality is symmetric, and on inequality the guard
clears both operands.  Therefore its surviving registers and subsequent
source control are identical to PA194H.  The change removes three source
sites.  Its additional physical fold is $41\mapsto330$, with actions
$0L31$ and $1L46$; all four local history/read-bit cases are rechecked in
the final machine together with the seven inherited AX folds.

The literal constructors produce PA194H with 3653 sites and 4754 source
outcomes, and PA194HS with 3650 sites and 4750 outcomes.  Their entire
resolved source graphs and complete tables are delivered.  Each outcome
has a universal symbolic physical contract over arbitrary natural
registers, both decrement branches, arbitrary affine header gap and an
arbitrary nonempty finite spare list.  The independent verifier derives
the four register-list lemmas again from the literal rows.  Exact blank
initialization reaches the clean source site 25 with logical file
$(2,0,\ldots,0)$ after 57922 and 57768 transitions respectively; the
complete finite-support tape is compared, including every spare block.

Positive physical durations give both directions of finite-trace
equivalence.  A source halt is simulated by concatenating finite positive
segments.  Conversely, if the physical run first halted at time $T$ while
the source had not halted, the first $T+1$ source steps would produce a
longer nonhalting physical trace, contradicting determinism.  This finite
argument and the loop certificates are formalizable in PA.  Consequently,
for either literal table $M$,
\[
 \mathrm{PA}\vdash
 \bigl(\neg\operatorname{Halt}(M,\mathrm{blank})
             \longleftrightarrow\operatorname{Con}(\mathrm{PA})\bigr).
\]

\paragraph{Why this is a different command design.}
The original no-push PA194N halts on Shift from $(1,0,0,0)$, whereas the
hybrid clears that store without halting.  The original cached-result
PA194C sends $(0,w,0,20)$ under $\mathsf{MP}_2$ to $(0,w,0,2)$, whereas
the hybrid sends it to $(2,w,2,0)$, since $I(0,2)=20$.  The new proof
compiler accounts for both differences.  These changes affect the specified
command behavior, rather than merely renaming states, tape directions or
formula parameters.
